%%%%%%%%%%%%%%%%%%%%%%%%%%%%%%%%%%%%%%%%%%%%%%%%%%%%%%%%%%%%%%%%%%%%%%
% Lớp 9 | Chương VIII. Xác suất của biến cố trong một số mô hình xác suất đơn giản | Xác suất của biến cố liên quan tới phép thử
% Dạng 9D8?26-2: Tính xác suất của biến cố
%%%%%%%%%%%%%%%%%%%%%%%%%%%%%%%%%%%%%%%%%%%%%%%%%%%%%%%%%%%%%%%%%%%%%%


% nguồn: 2026-11-15_09_de_tham_khao_quan_1_de_1_va_de_2 (nhập 15/11/2026)
\begin{bt}%[9D8H26-2]%[ma:47c49d11]
% nguon: Đề tham khảo Tuyển sinh 10 TP.HCM 2025-2026 – Quận 1 – Đề 1 – Bài 3
Một chiếc hộp chứa 3 viên bi gồm: một viên bi màu xanh (X), một viên bi màu vàng (V) và một viên bi màu đỏ (Đ). Các viên bi có cùng kích thước và khối lượng. Bạn An lần lượt lấy ra ngẫu nhiên từng viên bi từ trong hộp cho đến khi hộp hết viên bi.
\begin{enumerate}[a)]
	\item Xác định không gian mẫu của phép thử.
	\item Tính xác suất của mỗi biến cố sau:
	\begin{itemize}
		\item Biến cố $A$: ``Viên bi màu đỏ được lấy ra đầu tiên''.
		\item Biến cố $B$: ``Viên bi màu xanh được lấy ra trước viên bi màu vàng''.
		\item Biến cố $C$: ``Viên bi lấy ra lần cuối cùng không có màu xanh''.
	\end{itemize}
\end{enumerate}
\loigiai{
\begin{enumerate}[a)]
	\item Không gian mẫu $\Omega$ là tập hợp tất cả các thứ tự lấy ra 3 viên bi từ hộp:
	$\Omega = \{(\text{X}, \text{V}, \text{Đ}); (\text{X}, \text{Đ}, \text{V}); (\text{V}, \text{X}, \text{Đ}); (\text{V}, \text{Đ}, \text{X}); (\text{Đ}, \text{X}, \text{V}); (\text{Đ}, \text{V}, \text{X})\}$.\\
	Số phần tử của không gian mẫu là $n(\Omega) = 3! = 6$.
	\item Xác suất của từng biến cố:
	\begin{itemize}
		\item Các kết quả thuận lợi cho biến cố $A$ là: $(\text{Đ}, \text{X}, \text{V})$ và $(\text{Đ}, \text{V}, \text{X}) \Rightarrow n(A) = 2$.\\
		Do đó $P(A) = \dfrac{n(A)}{n(\Omega)} = \dfrac{2}{6} = \dfrac{1}{3}$.
		\item Các kết quả thuận lợi cho biến cố $B$ (bi màu xanh xuất hiện trước bi màu vàng) là: $(\text{X}, \text{V}, \text{Đ})$, $(\text{X}, \text{Đ}, \text{V})$ và $(\text{Đ}, \text{X}, \text{V}) \Rightarrow n(B) = 3$.\\
		Do đó $P(B) = \dfrac{n(B)}{n(\Omega)} = \dfrac{3}{6} = \dfrac{1}{2}$.
		\item Biến cố $C$ xảy ra khi viên bi cuối cùng là màu vàng hoặc màu đỏ. Các kết quả thuận lợi cho biến cố $C$ là: $(\text{X}, \text{V}, \text{Đ})$, $(\text{X}, \text{Đ}, \text{V})$, $(\text{V}, \text{X}, \text{Đ})$ và $(\text{Đ}, \text{X}, \text{V}) \Rightarrow n(C) = 4$.\\
		Do đó $P(C) = \dfrac{n(C)}{n(\Omega)} = \dfrac{4}{6} = \dfrac{2}{3}$.
	\end{itemize}
\end{enumerate}
}
\end{bt}

% nguồn: 2026-11-15_09_de_tham_khao_quan_1_de_1_va_de_2 (nhập 15/11/2026)
\begin{bt}%[9D8H26-2]%[ma:d360d887]
% nguon: Đề tham khảo Tuyển sinh 10 TP.HCM 2025-2026 – Quận 1 – Đề 2 – Bài 3
Một chiếc hộp có chứa $5$ tấm thẻ cùng loại, được đánh số lần lượt là $3; 5; 6; 7; 9$. Lấy ra ngẫu nhiên đồng thời $2$ tấm thẻ từ hộp.
\begin{enumerate}[a)]
	\item Xác định không gian mẫu và số kết quả có thể xảy ra của phép thử.
	\item Tính xác suất của biến cố $A$: ``Tích các số ghi trên $2$ tấm thẻ chia hết cho $3$''.
\end{enumerate}
\loigiai{
\begin{enumerate}[a)]
	\item Không gian mẫu $\Omega$ gồm tất cả các tập hợp con gồm $2$ thẻ chọn từ $5$ thẻ:
	$\Omega = \{(3; 5); (3; 6); (3; 7); (3; 9); (5; 6); (5; 7); (5; 9); (6; 7); (6; 9); (7; 9)\}$.\\
	Số kết quả có thể xảy ra của phép thử là $n(\Omega) = 10$.
	\item Tích của hai số chia hết cho $3$ khi và chỉ khi ít nhất một trong hai thẻ mang số chia hết cho $3$.\\
	Trong $5$ thẻ đã cho, các số chia hết cho $3$ là: $3; 6; 9$ (gồm $3$ thẻ); các số không chia hết cho $3$ là: $5; 7$ (gồm $2$ thẻ).\\
	Cặp duy nhất có tích không chia hết cho $3$ là $(5; 7)$ (vì $5\cdot 7 = 35$).\\
	Do đó số kết quả thuận lợi cho biến cố $A$ là: $n(A) = 10 - 1 = 9$.\\
	Xác suất của biến cố $A$ là:
	$P(A) = \dfrac{n(A)}{n(\Omega)} = \dfrac{9}{10} = 0{,}9$.
\end{enumerate}
}
\end{bt}

% nguồn: 2026-11-15_09_de_tham_khao_quan_3_de_1_va_de_2 (nhập 15/11/2026)
\begin{bt}%[9D8H26-2]%[ma:792d9d82]
% nguon: Đề tham khảo Tuyển sinh 10 TP.HCM 2025-2026 – Quận 3 – Đề 2 – Bài 3
Nhân dịp tổng kết năm học, lớp 9A có $2$ bạn nam hát hay là Tuấn, Hùng và $2$ bạn nữ hát hay là Lan, Hồng. Giáo viên chủ nhiệm chọn ngẫu nhiên $2$ bạn để hát song ca.
\begin{enumerate}[a)]
	\item Hãy liệt kê tất cả các cách chọn ngẫu nhiên $2$ bạn để hát song ca.
	\item Tính xác suất của mỗi biến cố sau:
	\begin{itemize}
		\item Biến cố $A$: ``Trong $2$ bạn được chọn có đúng $1$ bạn nam và $1$ bạn nữ''.
		\item Biến cố $B$: ``Trong $2$ bạn được chọn có bạn Lan''.
	\end{itemize}
\end{enumerate}
\loigiai{
\begin{enumerate}[a)]
	\item Tập hợp tất cả các cặp gồm $2$ bạn được chọn từ $4$ bạn (không phân biệt thứ tự) là không gian mẫu:
	$\Omega = \{(\text{Tuấn}, \text{Hùng}); (\text{Tuấn}, \text{Lan}); (\text{Tuấn}, \text{Hồng}); (\text{Hùng}, \text{Lan}); (\text{Hùng}, \text{Hồng}); (\text{Lan}, \text{Hồng})\}$.\\
	Số phần tử của không gian mẫu là $n(\Omega) = 6$.
	\item Xác suất của từng biến cố:
	\begin{itemize}
		\item Các kết quả thuận lợi cho biến cố $A$ (gồm 1 nam và 1 nữ) là:
		$(\text{Tuấn}, \text{Lan})$, $(\text{Tuấn}, \text{Hồng})$, $(\text{Hùng}, \text{Lan})$ và $(\text{Hùng}, \text{Hồng}) \Rightarrow n(A) = 4$.\\
		Do đó xác suất của biến cố $A$ là:
		$P(A) = \dfrac{n(A)}{n(\Omega)} = \dfrac{4}{6} = \dfrac{2}{3}$.
		\item Các kết quả thuận lợi cho biến cố $B$ (có bạn Lan) là:
		$(\text{Tuấn}, \text{Lan})$, $(\text{Hùng}, \text{Lan})$ và $(\text{Lan}, \text{Hồng}) \Rightarrow n(B) = 3$.\\
		Do đó xác suất của biến cố $B$ là:
		$P(B) = \dfrac{n(B)}{n(\Omega)} = \dfrac{3}{6} = \dfrac{1}{2}$.
	\end{itemize}
\end{enumerate}
}
\end{bt}
